\section{Abstract}

This thesis studies price tracking in a synthetic, zero-fee constant-product automated market maker (AMM) against an exogenous centralized-exchange (CEX) reference. It separates the swap rule's algebraic implications from effects that depend on agent behavior. In an isolated pool, an arbitrageur limited to a fraction $c$ of the input reserve can reduce out-of-band log error by at most $2\log(1+c)$ per tick. For a log shock of $0.20$ and a band of $0.003$, the idealized time to the band is 20 ticks at $c=0.005$, four at $c=0.03$, and two at $c=0.10$. This calculation explains the direction of the configured arbitrage result; simulation measures departures under noise and other rules.

Across 300 Monte Carlo seeds and 13,500 runs, the high-capacity scenario has a mean excess-error area of $0.10$, compared with $2.36$ without arbitrage. The result is conditional on the specified rule and parameter grid, and the sign of the effect is mechanically built into that rule. The current withdrawal trigger fires in every E3 run, usually at the first tick. A post-hoc maximum absolute tracking-error measure is higher under E3 only at depth factor $0.5$ with fixed absolute orders (paired mean differences $0.0104$--$0.0519$ across shock cells); the contrast is numerically zero at depth factors 1 and 2. Withdrawal also reduces the reserve-linked arbitrage cap, so this secondary result does not isolate a withdrawal effect or establish general liquidity amplification.

The E3 persistence result is simulator-specific: with the nine design cells equally weighted, 31.2\,\% of recovered paths later left the no-arbitrage band; at depth factor $0.5$ the conditional share was 98.8\,\%, while no later departures occurred at depth factors 1 and 2. This motivates recording first return and later departure separately in future event-replay studies, but it is not a real-market rate. The two-state reduced map improves on persistence in some conditions but does not outperform a scalar AR(1). A band-aware map improves recursive replay error by about 1.7\,\%, concentrated in one condition and using the realized future CEX path; it is not a real-time forecasting result. The contribution is a reproducible, mechanism-focused simulation study that distinguishes algebraic results from rule-dependent sensitivities and reports where its reduced models fail. Event replay with synchronized market data and realistic costs remains future work.

\textbf{Keywords:} automated market maker; arbitrage; price tracking; liquidity withdrawal; agent-based simulation; Monte Carlo experiment
